\scanpage{23}
\mspage{6)}
\[
(9)\quad\diagram{\begin{tikzcd}
&1\arrow[d]&1\arrow[d]&&\\
&\Pi_{X,x}\arrow[r,equal]\arrow[d]&\Pi_{X,x}\arrow[d]&&\\
1\arrow[r,no head]&\mathcal G_{X,x}\arrow[r,no head]\arrow[d,no head]&\mathcal M_{X,x\,K}\arrow[r,no head]\arrow[d]&\Gamma_{\overline K/K}\arrow[r,no head]\arrow[d,equal]&1\\
1\arrow[r,no head]&\Gamma_X\arrow[r,no head]\arrow[d]&\mathcal N_{X,x\,K}\arrow[r,no head]\arrow[d]&\Gamma_{\overline K/K}\arrow[r,no head]&1\\
&1&1&&
\end{tikzcd}}
\]
correspondant à
\[
(10)\qquad\overbrace{\underbrace{\Pi_{X,x}\subset\mathcal G_{X,x}}_{\Gamma_X}
\underbrace{{}\subset\mathcal M_{X,x,K}}_{\Gamma_{\overline K/K}}}^{\mathcal N_{X,K}}
\]
(\uline{NB} On n'a pas encore utilisé que $(X,x)$ soit défini sur $K$, i.e. que
$x$ \inserted{$(\in\mathcal X_{g,\nu}(k))$} provienne d'un pt à val. dans $K$\ldots).
À vrai dire, la situation (9) \inserted{ou (10)} se déduit de la situation
``absolue'' (i.e. sur $\mathbb Q$) (1), par \unsure{la|le} image inverse de structures
d'extensions (de $\Gamma_{\widetilde{\mathbb Q}/\mathbb Q}$ par
$\mathcal G_{X,x}\supset\Pi_{X,x}$) à l'aide de l'hom. de groupes
\[
(11)\qquad\Gamma_{\overline K/K}\longrightarrow\Gamma_{\widetilde{\mathbb Q}/\mathbb Q}
\]
\mspage{22}
